%% LyX 2.4.2.1 created this file.  For more info, see https://www.lyx.org/.
%% Do not edit unless you really know what you are doing.
\documentclass[11pt,american,english]{article}
\usepackage[utf8]{inputenc}
\usepackage{xcolor}
\usepackage{babel}
\usepackage{mathtools}
\usepackage{dsfont}
\usepackage{amsmath}
\usepackage{amsthm}
\usepackage{amssymb}
\usepackage[letterpaper]{geometry}
\geometry{verbose}
\usepackage[authoryear]{natbib}
\usepackage[bookmarks=false,
 breaklinks=false,pdfborder={0 0 1},backref=section,colorlinks=true]
 {hyperref}
\hypersetup{
 naturalnames=true,linkcolor=blue,citecolor=blue,urlcolor=blue}

\makeatletter
%%%%%%%%%%%%%%%%%%%%%%%%%%%%%% User specified LaTeX commands.
% AER-Article.tex for AEA last revised 22 June 2011


% General packages from this template
\usepackage[english]{babel}
\usepackage{graphicx}
\usepackage{amsfonts}
\usepackage{amsthm}
\usepackage{caption}
\usepackage{subcaption}
\usepackage{enumerate}
\usepackage[labelfont=sc]{caption}
\usepackage[capposition=top]{floatrow}
\usepackage{float}



% New packages for this article
\usepackage{siunitx}
\usepackage{pdflscape}
\usepackage{setspace}


% General setup from the template
\setlength{\abovecaptionskip}{0pt}
\setlength{\belowcaptionskip}{0 pt}

\newcommand{\fnote}[1]{\captionsetup{font=footnotesize}\caption*{#1}}

\linespread{1.3}

\newtheorem{theorem}{\textit{Theorem}}\newtheorem{proposition}[theorem]{\textit{Proposition}}\newtheorem{lemma}{\textit{Lemma}}\newtheorem{corollary}{\textit{Corollary}}\newtheorem{definition}{\textit{Definition}}\renewcommand{\qedsymbol}{$\blacksquare$}
% \newenvironment{definition}[1][Definition.]{\begin{trivlist}
% \item[\hskip \labelsep {\bfseries #1}]}{\end{trivlist}}
\newenvironment{assumption}[1][Assumption.]{\begin{trivlist}
\item[\hskip \labelsep {\bfseries #1}]}{\end{trivlist}}
\newenvironment{remark}[1][Remark]{\begin{trivlist}
\item[\hskip \labelsep {\bfseries #1}]}{\end{trivlist}}

\renewcommand*{\theenumi}{\roman{enumi}}
\renewcommand*{\labelenumi}{(\theenumi)}


\renewcommand{\@biblabel}[1]{}


\renewcommand{\thetable}{\Roman{table}}
\captionsetup[table]{labelformat=simple, labelsep=newline, textfont = {sc}}
\captionsetup[figure]{labelformat=simple}



% New setup for this article

\newcommand{\MatrixVector}[1]{{\mbox{\boldmath $#1$}}}

\DeclareMathOperator{\sign}{sign}
\DeclareMathOperator{\Var}{Var}
\DeclareMathOperator{\ForAll}{\,\forall\,}

\renewcommand{\textfraction}{0.01}
\renewcommand{\topfraction}{0.99}
\renewcommand{\bottomfraction}{0.99}

\newcommand*{\diff}{\mathop{}\!\mathrm{d}}
\newcommand*{\Diff}[1]{\mathop{}\!\mathrm{d}#1}
\newcommand{\ME}[2]{\ensuremath{#1^{\mathrm{me}}_{#2}}}
\newcommand{\Central}{^{\mathrm{c}}}
\newcommand{\Decentral}{^{\mathrm{d}}}
\newcommand{\Unemployed}{^{\mathrm{u}}}
\newcommand{\Employed}{^{\mathrm{e}}}
\newcommand{\Out}{^{\mathrm{o}}}
\newcommand{\Entre}{^{\mathrm{e}}}
\newcommand{\InvestShock}{^{\mathrm{I}}}
\newcommand{\House}{^{\mathrm{h}}}
\newcommand{\Transit}{^{\mathrm{ue}}}
\newcommand{\search}{^{\mathrm{s}}}
\newcommand{\Search}{_{\mathrm{s}}}
\newcommand{\Consume}{_{\mathrm{c}}}
\newcommand{\SocialPlanner}{_{\mathrm{sp}}}
\newcommand{\Retire}{^{\mathrm{r}}}
\newcommand{\SteadyState}{_{\!\mathrm{ss}}}
\newcommand{\Labor}{^{\mathrm{y}}}
\newcommand{\labor}{_{\mathrm{y}}}
\newcommand{\Capital}{^{\mathrm{k}}}


\usepackage{tikz}% for timeline
\usepackage{rotating}
\usetikzlibrary{arrows}
\usetikzlibrary{shapes,snakes}
\usetikzlibrary{positioning} 

\usetikzlibrary{calc}
\usepackage{relsize}


%\usepackage{amsmath, amssymb,mathrsfs}  
\usepackage{siunitx}
\sisetup{
input-symbols = {()},
group-digits  = false,
explicit-sign
}

\frenchspacing

% \usepackage[none]{hyphenat}
% \usepackage{microtype}
% \DisableLigatures{encoding = *, family = * }

\providecommand{\noopsort}[1]{}
\renewcommand{\topfraction}{0.9} % max fraction of floats at top
\renewcommand{\bottomfraction}{0.8} % max fraction of floats at bottom
\setcounter{topnumber}{2}
\setcounter{bottomnumber}{2}
\setcounter{totalnumber}{4} % 2 may work better
\setcounter{dbltopnumber}{2} % for 2-column pages
\renewcommand{\dbltopfraction}{0.9} % fit big float above 2-col. text
\renewcommand{\textfraction}{0.07} % allow minimal text w. figs
%   Parameters for FLOAT pages (not text pages):
\renewcommand{\floatpagefraction}{0.7}	% require fuller float pages
    % N.B.: floatpagefraction MUST be less than topfraction !!
\renewcommand{\dblfloatpagefraction}{0.7}	% require fuller float pages

\usepackage{tabularx}
\usepackage{booktabs}
\usepackage{textpos}% for exact text positions
\usepackage{graphicx}
\usepackage{epstopdf}


% code for epigraph

\newenvironment{chapquote}[2][2em]{\setlength{\@tempdima}{#1}%
   \def\chapquote@author{#2}%
   \parshape 1 \@tempdima \dimexpr\textwidth-2\@tempdima\relax%
   \itshape}{\par\normalfont\hfill--\ \chapquote@author\hspace*{\@tempdima}\par\bigskip}

\makeatother

\begin{document}
\title{{\Large\textbf{Model: Firm Dynamics and Political Capital}}}
\date{{\small\today}}

\maketitle
 

\section*{Households}

\paragraph{Preferences.}

A representative household has preferences 
\[
\sum_{t=0}^{\infty}\beta^{t}\log(c_{t}),\qquad\beta\in(0,1).
\]

Aggregate consumption is Cobb--Douglas between the non-differentiated
good $c_{t}^{n}$ and the political composite $c_{t}^{p}$: 
\[
c_{t}=(c_{t}^{n})^{1-\xi}(c_{t}^{p})^{\xi},\qquad\xi\in(0,1).
\]

Political consumption is a linear aggregator of $c_{t}^{p}(i)\geq0$
over a time-varying set of varieties $\mathcal{I}_{t}$ with weights
$\omega_{it}>0$: 
\[
c_{t}^{p}=\int_{\mathcal{I}_{t}}\omega_{it}\,c_{t}^{p}(i)\,di.
\]
Let $p_{t}(i)$ denote the real price of variety $i$, and let $p_{t}^{p}$
denote the real shadow price per effective unit of $c_{t}^{p}$. The
taste shifter $\omega_{it}$ follows a persistent AR(1) process in
logs: 
\[
\log\omega_{i,t+1}=\rho_{\omega}\log\omega_{it}+\sigma_{\omega}\varepsilon_{i,t+1}^{\omega},\qquad\rho_{\omega}\in(0,1),
\]
with $\varepsilon_{i,t+1}^{\omega}\sim i.i.d.\ N(0,1)$.

\paragraph{Optimization.}

Households trade a one-period risk-free bond that pays $1$ unit of
the numeraire at $t+1$. Let $R_{t}$ denote the gross real interest
rate and let $b_{t+1}$ denote bond holdings chosen at $t$. Let $\Pi_{t}$
denote lump-sum transfers from firms.

The budget constraint is 
\[
c_{t}^{n}+\int_{\mathcal{I}_{t}}p_{t}(i)c_{t}^{p}(i)\,di+\frac{b_{t+1}}{R_{t}}=\Pi_{t}+b_{t}.
\]
Define total real consumption expenditure 
\[
C_{t}\equiv c_{t}^{n}+\int_{\mathcal{I}_{t}}p_{t}(i)c_{t}^{p}(i)\,di=c_{t}^{n}+p_{t}^{p}c_{t}^{p}.
\]


\paragraph{Key conditions.}

(i) Political varieties: 
\[
\frac{p_{t}(i)}{\omega_{it}}\ge p_{t}^{p}\ \ \forall i\in\mathcal{I}_{t},\qquad\frac{p_{t}(i)}{\omega_{it}}=p_{t}^{p}\ \ \text{if }c_{t}^{p}(i)>0.
\]

(ii) Non-differentiated vs.\ political composite: 
\[
p_{t}^{p}=\frac{\xi}{1-\xi}\frac{c_{t}^{n}}{c_{t}^{p}},\qquad c_{t}^{n}=(1-\xi)C_{t},\quad p_{t}^{p}c_{t}^{p}=\xi C_{t}.
\]

(iii) Risk-free bond Euler equation: 
\[
\frac{1}{C_{t}}=\beta R_{t}\,\mathbb{E}_{t}\!\left[\frac{1}{C_{t+1}}\right].
\]


\section*{Firms}

Firms are indexed by $j\in[0,1]$.

\paragraph{Technologies.}

Firms have access to a technology to produce the undifferentiated
good given by: 
\[
y_{jt}=z_{jt}k_{jt}^{\alpha},\qquad\alpha\in(0,1),
\]
with productivity follows a persistent AR(1) process (in logs): 
\[
\log z_{j,t+1}=(1-\rho_{z})\log\bar{z}+\rho_{z}\log z_{jt}+\sigma_{z}\varepsilon_{j,t+1}^{z},\qquad\rho_{z}\in(0,1),
\]
with $\varepsilon_{j,t+1}^{z}\sim i.i.d.\ N(0,1)$.

Firms have access to a technology to accumulate physical capital.
\[
k_{j,t+1}=(1-\delta)k_{jt}+i_{jt},\qquad\delta\in(0,1).
\]

Firms can differentiate their output by accumulating political capital
$x_{jt}\in\{0,1\}$. Firms that currently have no political capital
can accumulated it with political engagement $e_{jt}\in\{0,1\}$ (that
maps to political speech). Political capital depreciates with probability
$\delta_{x}$. Let $\nu_{j,t+1}\in\{0,1\}$ denote a political-capital
survival indicator: $\nu_{j,t+1}=1$ if the firm retains its political
capital between $t$ and $t+1$, and $\nu_{j,t+1}=0$ otherwise, with
\[
\Pr(\nu_{j,t+1}=1)=1-\delta_{x},\qquad\Pr(\nu_{j,t+1}=0)=\delta_{x}.
\]
Political-capital then follows the law of motion: 
\[
x_{j,t+1}=x_{jt}\nu_{j,t+1}+(1-x_{jt})e_{jt}.
\]

Every period they exit with probability $\chi$ after production and
are replaced by an equal mass of firms. When firms exit they transfer
their net worth to households.

\paragraph{Prices.}

Let $\iota:[0,1]\to[0,1]$ be a one-to-one mapping from firm indices
$j$ to potential political varieties $i=\iota(j)$. Define the variety-specific
taste shifter faced by firm $j$ as 
\[
\omega_{jt}\equiv\omega_{\iota(j),t}.
\]

Let $p_{t}^{p}$ denote the competitive price per effective unit of
the political composite. The firm takes prices as given and receives
the unit output price 
\[
p_{t}(x_{jt},\omega_{jt})=\begin{cases}
1, & x_{jt}=0,\\
\omega_{jt}\,p_{t}^{p}, & x_{jt}=1.
\end{cases}
\]


\paragraph{Financing.}

The firm can issue one-period debt. Let $d_{jt}$ be debt due at $t$
(promised repayment made at $t-1$) and let $d_{j,t+1}$ be promised
repayment next period. The gross real interest rate is $R_{t}$.

The flow-of-funds constraint is given by: 
\[
\mathrm{div}_{jt}=\pi_{t}(k_{jt},x_{jt},z_{jt},\omega_{jt},e_{jt})-i_{jt}-\Phi(k_{j,t+1},k_{jt})+\frac{d_{j,t+1}}{R_{t}}-d_{jt}.
\]
where $\pi_{t}(k_{jt},x_{jt},z_{jt},\omega_{jt},e_{jt})=p_{t}(x_{jt},\omega_{jt})\,z_{jt}k_{jt}^{\alpha}-c_{x}e_{jt}$
denote operating earnings and $\Phi(k_{j,t+1},k_{jt})$ denote a capital
adjustment cost in units of the undifferentiated good.

Firms face an earning-based borrowing constraint of the form: 
\[
d_{j,t+1}\le\phi\,\pi_{t}(k_{jt},x_{jt},z_{jt},\omega_{jt},e_{jt}),\qquad\phi\ge0.
\]
In addition, they cannot issue equity $\mathrm{div}_{jt}\ge0$.

\paragraph{Optimization.}

Let the firm’s idiosyncratic state be $s=(k,x,d,z,\omega)$ and $\omega\equiv\omega_{\iota(j)}$
the variety if the firm has political capital. Let $\Lambda_{t+1}$
denote the stochastic discount factor. Firms' recursive problem is
given by: 
\[
V_{t}(s)=\max_{\mathrm{div}\ge0,i,\,e\in\{0,1\},\,d'}\left\{ \mathrm{div}+\mathbb{E}_{t}\Big[\Lambda_{t+1}\Big((1-\chi)\,V_{t+1}(s')+\chi\,n_{t+1}(s')\Big)\Big]\right\} ,
\]
subject to 
\begin{align}
k' & =(1-\delta)k+i,\\
x' & =x\,\nu_{t+1}+(1-x)e,\\
\mathrm{div} & =\pi_{t}(k,x,z,\omega,e)-i-\Phi(k',k)+\frac{d'}{R_{t}}-d,\\
\pi_{t}(k,x,z,\omega,e) & =p_{t}(x,\omega)\,z\,k^{\alpha}-c_{x}e,\\
d' & \le\phi\,\pi_{t}(k,x,z,\omega,e),
\end{align}
where net worth is defined by $n_{t}(s)\equiv p_{t}(x,\omega)\,z\,k^{\alpha}-c_{x}\hat{e}_{t}(s)+(1-\delta)k-d$

\paragraph{Political engagement decision.}

For a firm with $x=0$, political engagement $e\in\{0,1\}$ pins down
next-period political status: $x'=e$. Define continuation values
conditional on the engagement choice: 
\[
\mathcal{V}_{t}^{n}(k,d,z,\omega)\equiv\max_{i,\,d'}\left\{ \mathrm{div}\big|_{e=0}+\mathbb{E}_{t}\!\left[\Lambda_{t+1}\Big((1-\chi)\,V_{t+1}(k',0,d',z',\omega')+\chi\,n'\Big)\right]\right\} ,
\]
\[
\mathcal{V}_{t}^{p}(k,d,z,\omega)\equiv\max_{i,\,d'}\left\{ \mathrm{div}\big|_{e=1}+\mathbb{E}_{t}\!\left[\Lambda_{t+1}\Big((1-\chi)\,V_{t+1}(k',1,d',z',\omega')+\chi\,n'\Big)\right]\right\} .
\]
The optimal engagement rule satisfies 
\[
e=1\quad\Longleftrightarrow\quad\mathcal{V}_{t}^{p}(k,d,z,\omega)\ \ge\ \mathcal{V}_{t}^{n}(k,d,z,\omega).
\]
For $x'=1$, the firm receives the political-market price $p_{t}(1,\omega)=\omega\,p_{t}^{p}$,
which can exceed $1$, but pays the engagement cost $c_{x}$ and becomes
exposed to additional cash-flow volatility through $\omega$.

\section*{Equilibrium}

\paragraph{Set of political varieties.}

Let the set of politically active firms at time $t$ be 
\[
\mathcal{J}_{t}^{p}\equiv\{j\in[0,1]:x_{jt}=1\},\qquad X_{t}\equiv\int_{0}^{1}x_{jt}\,dj.
\]
The set of political varieties available for trade at time $t$ is
endogenous and given by 
\[
\mathcal{I}_{t}\equiv\{\,\iota(j):j\in\mathcal{J}_{t}^{p}\,\},
\]
with measure 
\[
\int_{\mathcal{I}_{t}}di=X_{t}.
\]


\paragraph{Market clearing.}

For each active variety $i\in\mathcal{I}_{t}$, let $j(i)$ denote
the inverse mapping (so that $\iota(j(i))=i$). Variety-by-variety
clearing: 
\[
c_{t}^{p}(i)=y_{j(i),t}=z_{j(i),t}k_{j(i),t}^{\alpha},\qquad\forall i\in\mathcal{I}_{t}.
\]
Aggregate political composite: 
\[
c_{t}^{p}=\int_{\mathcal{I}_{t}}\omega_{it}\,c_{t}^{p}(i)\,di=\int_{0}^{1}x_{jt}\,\omega_{\iota(j),t}\,z_{jt}k_{jt}^{\alpha}\,dj.
\]

The non-differentiated good is used for household consumption, physical
investment, capital adjustment costs, and political engagement costs:
\[
c_{t}^{n}+\int_{0}^{1}i_{jt}\,dj+\int_{0}^{1}\Phi(k_{j,t+1},k_{jt})\,dj+c_{x}\int_{0}^{1}e_{jt}\,dj=\int_{0}^{1}(1-x_{jt})\,z_{jt}k_{jt}^{\alpha}\,dj.
\]

Bond market clearing: 
\[
b_{t+1}=\int_{0}^{1}d_{j,t+1}\,dj.
\]


\section*{Solving the firm's problem}

The value of a continuing firm, conditional on its engagement choice
$e$, is given by 
\begin{align*}
v_{t}^{\text{cont}}\left(s\,|\,e\right)= & \max_{k^{\prime},d^{\prime}}\,n_{t}\left(s\right)-k^{\prime}-\Phi\left(k^{\prime},k\right)+\frac{d^{\prime}}{R_{t}}+\Lambda_{t,t+1}\mathbb{E}_{t}\left[\left(1-\chi\right)v\left(s^{\prime}\right)+\chi n^{\prime}\,|\,z,\omega\right],
\end{align*}
subject to $\text{div}_{t}\geq0$ and $d^{\prime}\le\phi\,\pi_{t}\left(k,x,z,\omega,e\right)$,
where $n_{t}\left(s\right)\equiv p_{t}\left(x,\omega\right)\,z\,k^{\alpha}-c_{x}e+\left(1-\delta\right)k-d$
is net worth, the expectation $\mathbb{E}_{t}$ is taken over the
distribution of next period's exogenous idiosyncratic state $\left(z^{\prime},\omega^{\prime}\right)$,
and $v\left(s\right)=\max\left\{ v_{t}^{\text{cont}}\left(s\,|\,e=0\right),v_{t}^{\text{cont}}\left(s\,|\,e=1\right)\right\} $.

Writing down the Lagrangian for this problem, we have 
\begin{align*}
\mathcal{L} & =\left(1+\lambda_{t}\left(s\right)\right)\left(n_{t}\left(s\right)-k^{\prime}-\Phi\left(k^{\prime},k\right)+\frac{d^{\prime}}{R_{t}}\right)+\Lambda_{t,t+1}\mathbb{E}_{t}\left[\left(1-\chi\right)v\left(s^{\prime}\right)+\chi n^{\prime}\,|\,z,\omega\right]\\
 & +\xi_{t}\left(s\right)\left(\phi\,p_{t}\left(x,\omega\right)zk^{\alpha}-c_{x}e-d^{\prime}\right),
\end{align*}

First order condition for capital: 
\begin{align*}
\left(1+\lambda_{t}\left(s\right)\right)\left(1+\Phi_{1}\left(k^{\prime},k\right)\right) & =\Lambda_{t,t+1}\mathbb{E}_{t}\left[\left(1-\chi\right)\frac{\partial v\left(s^{\prime}\right)}{\partial k^{\prime}}+\chi\frac{\partial n^{\prime}}{k^{\prime}}\right],
\end{align*}

The envelope condition for capital is given by 
\begin{align*}
\frac{\partial v\left(s\right)}{\partial k} & =\left(1+\lambda_{t}\left(s\right)\right)\left(p_{t}(x,\omega)z\alpha k^{\alpha-1}+(1-\delta)-\Phi_{2}\left(k^{\prime},k\right)\right)+\xi_{t}\left(s\right)\phi p_{t}(x,\omega)z\alpha k^{\alpha-1}.
\end{align*}

Using these two conditions, and defining $\text{MPK}_{t+1}\equiv p_{t+1}(x^{\prime},\omega^{\prime})z^{\prime}\alpha(k^{\prime})^{\alpha-1}$,
we can rewrite the first order condition for capital as, 
\begin{align}
\left(1+\lambda_{t}\left(s\right)\right)\left(1+\Phi_{1}\left(k^{\prime},k\right)\right) & =\Lambda_{t,t+1}\mathbb{E}_{t}\left[\left(1+\left(1-\chi\right)\lambda(s^{\prime})\right)\left(\text{MPK}_{t+1}+1-\delta\right)\right.\label{eq:foc_k}\\
 & -\left.\left(1-\chi\right)(1+\lambda\left(s^{\prime}\right))\Phi_{2}\left(k^{\prime\prime},k^{\prime}\right)+\left(1-\chi\right)\xi\left(s^{\prime}\right)\phi\text{MPK}_{t+1}\right]\nonumber 
\end{align}

First order condition for debt: 
\begin{align*}
\left(1+\lambda_{t}\left(s\right)\right)\frac{1}{R_{t}}+\mathbb{E}_{t}\left[\left(1-\chi\right)\frac{\partial v\left(s^{\prime}\right)}{\partial d^{\prime}}-\chi\right] & =\xi_{t}\left(s\right)
\end{align*}

The envelope condition for debt is given by 
\begin{align*}
\frac{\partial v\left(s\right)}{\partial d} & =-\left(1+\lambda_{t}\left(s\right)\right)
\end{align*}

Using these two conditions, we can rewrite the first order condition
for debt as, 
\begin{align}
\xi_{t}\left(s\right) & =\left(1+\lambda_{t}\left(s\right)\right)\frac{1}{R_{t}}-\Lambda_{t,t+1}\mathbb{E}_{t}\left[1+\left(1-\chi\right)\lambda\left(s^{\prime}\right)\right].\label{eq:d_foc}
\end{align}

We now have two first order conditions and the following complementary
slackness conditions: 
\begin{align*}
\lambda_{t}\left(s\right) & \geq0,\quad\text{div}_{t}\geq0,\quad\lambda_{t}\left(s\right)\text{div}_{t}=0,\\
\xi_{t}\left(s\right) & \geq0,\quad d^{\prime}\le\phi\,\pi_{t}\left(k,x,z,\omega,e\right),\quad\xi_{t}\left(s\right)\left(d^{\prime}-\phi\,\pi_{t}\left(k,x,z,\omega,e\right)\right)=0.
\end{align*}


\section*{Partitioning the state space}

Our state space is $s=\left(k,x,d,z,\omega\right)$. Note that net
worth is linear in debt, so we can equivalently write the state space
as $s=(k,x,n,z,\omega)$. We can partition the state space into three
regions as a function of net worth $n\left(k,x,z,\omega\right)$.

\subsection*{Unconstrained firms}

A financial unconstrained firm choices are independent of its financial
position, can operate at its efficient scale, it is indifferent between
$d^{\prime}$ and $\text{div}$, and for which the marginal value
of savings in the same as for households. This means that $\lambda_{t+\tau}=\xi_{t+\tau}=0$
for $\tau\geq0$. To show that its choices are independent of its
financial position, we will guess and verify that these choices are
given by $k^{*}(k,x,z,\omega)$, and $d^{*}(k,x,z,\omega)$. To resolve
the indeterminacy between savings and dividends, we assume that these
firms implement a minimum savings policy, a firm's savings rule that
guarantees that the firm is never financially constrained in the future
under any possible paths of $(z,\omega)$. This policy is given by
the most debt (or least savings) that leave the firm unconstrained
in any future solves: 
\begin{equation}
d_{t}^{*}\left(k,x,z,\omega\right)=\min_{\left(z^{\prime},\omega^{\prime}\right)\in\mathbf{Z}\times\mathbf{\Omega}}\tilde{d}_{t+1}\left(k_{t}^{*}\left(k,x,z,\omega\right),x^{\prime},z^{\prime},\omega^{\prime}\right)\label{eq:msp}
\end{equation}
where 
\begin{align*}
\tilde{d}_{t}\left(k,x,z,\omega\right) & =\min_{e\in\left\{ 0,1\right\} }\tilde{\pi}_{t}\left(k,x,z,\omega,e\right)+\left(1-\delta\right)k-k^{*}\left(k,x,z,\omega\right)\\
 & -\Phi\left(k^{*}\left(k,x,z,\omega\right),k\right)+\frac{1}{R_{t+1}}d^{*}\left(k,x,z,\omega\right)
\end{align*}
This contraction guarantees that $\text{div}_{t+\tau}\geq0$ for $\tau\geq0$.
Note that we consider only the case in which there is engagement whenever
possible, as it would be the ``worst case scenario''. If a firm
can remain unconstrained in this case, it will remain unconstrained
in any case. This means that
\[
\text{arg}\min_{e\in\left\{ 0,1\right\} }\tilde{\pi}_{t}\left(k,x,z,\omega,e\right)=\begin{cases}
e=0 & \text{if }x=1\\
e=1 & \text{if }x=0
\end{cases}
\]

Define the problem of a firm adopting these policies as: 
\begin{align}
v_{t}^{*}\left(k,x,z,\omega\right) & =n_{t}-k^{*}\left(k,x,z,\omega\right)-\Phi\left(k^{*},k\right)+\frac{1}{R_{t}}d^{*}\left(k,x,z,\omega\right)\nonumber \\
 & +\Lambda_{t,t+1}\mathbb{E}_{t}\left[\left(1-\chi\right)v^{*}\left(k^{\prime},x^{\prime},z^{\prime},\omega^{\prime}\right)+\chi n^{\prime}\right]\label{eq:value_uncons}
\end{align}
Policies and continuation values are independent of $n$, thus we
can write $v^{*}\left(k,x,z,\omega\right)=n+\overline{v}^{*}\left(k,x,z,\omega\right)$,
where 
\begin{align*}
\overline{v}^{*}\left(k,x,z,\omega\right) & =-k^{*}\left(k,x,z,\omega\right)-\Phi\left(k^{*},k\right)+\frac{1}{R_{t}}d^{*}\left(k,x,z,\omega\right)\\
 & +\Lambda_{t,t+1}\mathbb{E}_{t}\left[\left(1-\chi\right)v^{*}\left(k^{\prime},x^{\prime},z^{\prime},\omega^{\prime}\right)+\chi n^{\prime}\right]
\end{align*}
Implying that the value of an unconstrained firm is linearly separable
in net worth (and that the marginal value of internal net worth is
1). From this simplified problem, we obtain the following optimality
condition for capital: 
\begin{align}
1+\Phi_{1}\left(k^{\prime},k\right) & =\Lambda_{t,t+1}\mathbb{E}_{t}\left[\text{MPK}_{t+1}+1-\delta-\left(1-\chi\right)\Phi_{2}\left(k^{\prime\prime},k^{\prime}\right)\right]\label{eq:k_efficient}
\end{align}
\foreignlanguage{american}{which is a frictionless (with-adjustment-costs)
policy. If these policies are feasible they are also optimal: once
the unconstrained allocation is feasible today and in any future time,
constraints do not distort choices, so the unconstrained allocation
is optimal. A firm with high-enough net worth that allows it to implement
these policies, is one in which the limited liability constraint is
not violated 
\[
n-k^{*}\left(k,x,z,\omega\right)-\Phi(k^{*},k)+\frac{1}{R_{t}}d^{*}\left(k,x,z,\omega\right)\geq0
\]
A condition which is satisfied only if 
\begin{equation}
n\geq\overline{n}\left(s\right)\equiv k^{*}\left(k,x,z,\omega\right)+\Phi(k^{*},k)-\frac{1}{R_{t}}d^{*}\left(k,x,z,\omega\right)\label{eq:cutoff_uncons}
\end{equation}
}

\subsection*{Constrained firms}

Constrained firms are defined as firms with $\lambda_{t}\left(s\right)>0$.
By construction, these firms cannot implement the unconstrained policies.
Among these firms, there are two subset of firms: (i) firms that are
currently constrained, so that the borrowing constraint is binding
and $\xi_{t}\left(s\right)>0$, and (ii) partially constrained firms,
which does not currently have a binding borrowing constraint, but
may be constrained in the future.

\subsubsection*{Currently constrained firms}

These firms have a binding borrowing constraint, so their debt choice
is pinned down by it 
\begin{equation}
d^{c}\left(s\right)=\phi\pi_{t}\left(k,x,z,\omega,e\right)\label{eq:dprime_currcons}
\end{equation}
Further, as their dividend constraint is binding, their capital choice
is determined by it 
\begin{equation}
n_{t}-k^{c}\left(s\right)-\Phi\left(k^{c}\left(s\right),k\right)+\frac{1}{R_{t}}d^{c}\left(s\right)=0\label{eq:kprime_currcorns}
\end{equation}


\subsubsection*{Partially constrained firms}

Finally, the partially constrained firms solve the system of optimality
conditions derived before. Denote their choices by $k^{p}\left(s\right)$,
and $d^{p}\left(s\right)$. These policies are feasible as long as
\[
n\geq\underline{n}\left(s\right)\equiv k^{p}\left(s\right)+\Phi(k^{p}\left(s\right),k)-\frac{1}{R_{t}}d^{p}\left(s\right)\geq0
\]

\selectlanguage{american}%

\section*{Stationary equilibrium: numerical solution}

\selectlanguage{english}%
Assume the following capital adjustment cost function: 
\begin{equation}
\Phi\left(k^{\prime},k\right)=\frac{\psi}{2}\left[\frac{k^{\prime}}{k}-\left(1-\delta\right)\right]^{2}k\label{eq:kadj_fnc}
\end{equation}
Given a vector of prices $\left\{ R_{t},\Lambda_{t,t+1},p_{t}^{p}\right\} $,
\foreignlanguage{american}{we can exploit the partition of the state
space defined before. First, let's solve the unconstrained firm problem.
We can then use their solution along with the optimality conditions
to solve the problem of all firms.}

\selectlanguage{american}%

\selectlanguage{english}%

\subsection*{Unconstrained firm decisions}
\begin{enumerate}
\item Guess a value $v^{*}\left(s\right)$, and a policy for $k^{*}\left(s\right)$.
Also, guess some engagement policy for these firms, $e^{*}\left(s\right)\in\left\{ 0,1\right\} $.
We do not need to guess $e^{*}\left(s\right)$ for firms $x=1$, as
we are assuming $e=0$ for them.
\item Determine the efficient scale of operation $k^{*}\left(s\right)$
using \eqref{eq:k_efficient}. Given the presence of convex adjustment
costs, we need to solve this numerically for $k^{*}\left(s\right)=k^{*}\left(k,x,{\color{lightgray}d},z,\omega\right)$.
Define: 
\begin{align*}
G\left(k^{\prime}|k,x,z,\omega,k^{*}\left(s\right)\right) & =1+\Phi_{1}\left(k^{\prime},k\right)\\
 & -\Lambda_{t,t+1}\mathbb{E}_{t}\left[p_{t+1}(x^{\prime},\omega^{\prime})z^{\prime}\alpha\left(k^{\prime}\right)^{\alpha-1}+1-\delta-\left(1-\chi\right)\Phi_{2}\left({\color{red}k^{*}}\mathinner{\color{red}\left(s^{\prime}\right)},k^{\prime}\right)\right]
\end{align*}
We can iterate over the policy function $k^{*}\left(s\right)$, while
solving $G\left(k^{\prime}\right)=0$, starting from an initial \textcolor{red}{guess}
for it. Do this until convergence. Note that everything is know from
the state space, except $x^{\prime}$, which we are taking from Step
1. 
\item Solve the minimum savings policy in \eqref{eq:msp}. 
\[
d_{t}^{*}\left(k,x,z,\omega\right)=\min_{\left(z^{\prime},\omega^{\prime}\right)\in\mathbf{Z}\times\mathbf{\Omega}}\tilde{d}_{t+1}\left(k_{t}^{*}\left(k,x,z,\omega\right),x^{\prime},z^{\prime},\omega^{\prime}\right)
\]
where 
\begin{align*}
\tilde{d}_{t}\left(k,x,z,\omega\right) & =\pi_{t}\left(k,x,z,\omega,e\right)+\left(1-\delta\right)k-k^{*}\left(k,x,z,\omega\right)\\
 & -\Phi\left(k^{*}\left(k,x,z,\omega\right),k\right)+\frac{1}{R_{t+1}}d^{*}\left(k,x,z,\omega\right)
\end{align*}
\item Update the value of the unconstrained firm \eqref{eq:value_uncons}
using the previous iteration of $k^{*}\left(s\right)$. When $x=0$,
do this for both engagement options $e\in\left\{ 0,1\right\} $, and
update $v^{*}\left(s\right)=\max\left\{ v^{*p}\left(s\right),v^{*n}\left(s\right)\right\} $,
and $e^{*}\left(s\right)=\text{argmax}_{e}\left\{ v^{*p}\left(s\right),v^{*n}\left(s\right)\right\} $. 

\textcolor{red}{Something to consider: Why do we want net worth to
include the cost of political engagement? It is considered as investment
in political capital for next period, $x^{\prime}$, similar to $k^{\prime}$.
It should show up on dividends, but not on net worth. Why is this
relevant: computational reasons: The continuation value when exiting
needs to be evaluated in $e^{\prime}$ as long as $e^{\prime}$ is
part of $\pi^{\prime}$ and hence of $n^{\prime}$ (given that exit
happens right after production, again, should we consider the cost
of $e^{\prime}$?)}
\item Compute the cutoff $\overline{n}$$\left(s\right)$ using \eqref{eq:cutoff_uncons}.
\end{enumerate}

\subsection*{Firm decisions}
\begin{enumerate}
\item Guess a value $v\left(s\right)$, a policy for $k\left(s\right)$
and $d\left(s\right)$, and a multipliers $\lambda\left(s\right)$
and $\xi\left(s\right)$. Also, guess some engagement policy firms,
$e\left(s\right)\in\left\{ 0,1\right\} $. 
\item For every point in the state space $s$, check if $n\left(s\right)\geq\overline{n}\left(s\right)$.
If so, set $k^{\prime}=k^{*}\left(s\right)$, $d^{\prime}=d^{*}\left(s\right)$,
$e=e^{*}\left(s\right)$, and $v=v^{*}\left(s\right)$ at those points
in $\mathbf{S}$. 
\item If not, let's solve the capital and debt policies for both values
$e=\left\{ 0,1\right\} $. Assume that the borrowing constraint is
not binding. Use the optimality condition for capital \eqref{eq:foc_k}
to update the $k^{\prime}$ policy using 
\begin{align*}
G\left(k^{\prime}\right) & =\left(1+\Phi_{1}\left(k^{\prime},k\right)\right)\left(1+\lambda_{t}\left(s\right)\right)\\
 & -\Lambda_{t,t+1}\mathbb{E}_{t}\left[\left(1+\left(1-\chi\right)\lambda(s^{\prime})\right)\left(p_{t+1}(x^{\prime},\omega^{\prime})z^{\prime}\alpha\left(k^{\prime}\right)^{\alpha-1}+1-\delta\right)\right.\\
 & \left.-\left(1-\chi\right)\Phi_{2}\left(k^{*}\left(s^{\prime}\right),k^{\prime}\right)\right]
\end{align*}
Iterate the policy function until convergence. 
\item Update $d\left(s\right)$ from the binding dividends constraint.
\item Update $\lambda\left(s\right)$ using \eqref{eq:d_foc}: 
\begin{align*}
\left(1+\lambda_{t}\left(s\right)\right) & =R_{t}\Lambda_{t,t+1}\mathbb{E}_{t}\left[1+\left(1-\chi\right)\lambda\left(s^{\prime}\right)\right]
\end{align*}
\item Check if borrowing constraint is actually not binding. If it is, compute
the debt policy using \eqref{eq:dprime_currcons}. Solve \eqref{eq:kprime_currcorns}
for the capital policy.
\item Compute the value 
\[
V_{t}\left(s\right)=\begin{cases}
\max\left\{ \mathcal{V}_{t}^{n}\left(k,d,z,\omega\right),\mathcal{V}_{t}^{p}\left(k,d,z,\omega\right)\right\}  & \text{if }x=0\\
\mathcal{V}_{t}^{n}\left(k,d,z,\omega\right) & \text{if }x=1
\end{cases}
\]
\end{enumerate}

\subsection*{Stationary distribution}

Let $\mu\left(k,x,d,z,\omega\right)$ denote the distribution of firms.
The distribution evolves according to
\begin{align*}
\mu_{t+1}\left(s^{\prime}\right) & =\left(1-\chi\right)\int_{\mathbf{K}\times\mathbf{D}\times\mathbf{Z}\times\mathbf{\Omega}}\mathds{1}\left\{ \left(k^{\prime}\left(s\right),x^{\prime}\left(s\right),d^{\prime}\left(s\right)\right)\right\} \mu\left(dk,0,dd,dz,d\omega\right)\\
 & +\left(1-\chi\right)\int_{\mathbf{K}\times\mathbf{D}\times\mathbf{Z}\times\mathbf{\Omega}}\left(1-\delta_{x}\right)\mathds{1}\left\{ \left(k^{\prime}\left(s\right),1,d^{\prime}\left(s\right)\right)\right\} \mu\left(dk,1,dd,dz,d\omega\right)\\
 & +\left(1-\chi\right)\int_{\mathbf{K}\times\mathbf{D}\times\mathbf{Z}\times\mathbf{\Omega}}\left(\delta_{x}\right)\mathds{1}\left\{ \left(k^{\prime}\left(s\right),0,d^{\prime}\left(s\right)\right)\right\} \mu\left(dk,1,dd,dz,d\omega\right)\\
 & +\chi\mu^{\text{new}}\left(s\right)
\end{align*}
Where $\mu^{\text{new}}\left(s\right)$ is the distribution of new
firms.
\end{document}
